\chapter{Outils d'automatisation et d'assistance à l'audit des SI}
\label{annexe:outils-audit}

\section*{Introduction}
\addcontentsline{toc}{section}{Introduction}

L'auditeur des systèmes d'information dispose aujourd'hui d'un écosystème
d'outils logiciels qui transforment profondément la nature et la portée de
ses missions. Ces outils ne remplacent pas le jugement professionnel~; ils
l'amplifient~: là où l'auditeur traditionnel sondait par échantillonnage,
l'outil d'analyse de données peut analyser la population entière des
transactions~; là où le questionnaire papier collectait des déclarations,
la plateforme GRC consolide des preuves documentées et traçables.

Cette annexe est organisée en quatre temps~: d'abord une grille de critères
de sélection indépendante des outils~; ensuite une classification fonctionnelle
qui regroupe les outils par famille plutôt que par éditeur~; puis des fiches
synthétiques par outil~; enfin un arbre de décision permettant de choisir
selon le contexte de la mission.

% ------------------------------------------------------------------------------
\section{Critères de sélection d'un outil d'audit}
\label{sec:criteres-outils}
% ------------------------------------------------------------------------------

Avant de comparer des outils, l'auditeur doit définir les critères qui
guideront son choix. Ces critères sont indépendants des outils~: ils
dépendent du contexte de la mission, du niveau de maturité de l'équipe
et des contraintes institutionnelles.

\subsection*{Critère~1. Périmètre fonctionnel couvert}

Un outil peut couvrir un ou plusieurs des cinq périmètres suivants~:
\begin{enumerate}
  \item \textbf{Analyse de données (AD)~:} extraction, manipulation,
        interrogation et visualisation de fichiers de transactions~;
  \item \textbf{Gestion de la mission (GM)~:} planification, papiers de
        travail, suivi des recommandations~;
  \item \textbf{GRC~:} gouvernance, gestion des risques et conformité,
        cartographie des risques, auto-évaluation des contrôles~;
  \item \textbf{Contrôle continu (CC)~:} surveillance automatisée des
        transactions et détection d'anomalies en temps réel ou quasi-réel~;
  \item \textbf{Extraction et transformation (ET)~:} lecture de formats
        de fichiers complexes (PDF, rapports propriétaires, exports ERP).
\end{enumerate}

\subsection*{Critère~2. Intégration aux référentiels}

Certains outils intègrent nativement les référentiels (COSO, COBIT, ISO,
IIA) comme cadres d'évaluation. Cette intégration réduit la charge de
paramétrage et renforce l'opposabilité des conclusions.

\subsection*{Critère~3. Courbe d'apprentissage et autonomie}

Un outil destiné à un auditeur seul ou à une petite équipe doit être
opérationnel rapidement~; un outil destiné à une institution d'audit
de grande taille peut justifier un investissement de formation plus long.
La distinction clé est~: \emph{outil pour auditeur individuel} (ACL
Desktop, IDEA) vs.\ \emph{plateforme institutionnelle} (SAP GRC, Oracle GRC).

\subsection*{Critère~4. Accessibilité et coût}

\begin{itemize}
  \item \textbf{Outils commerciaux sur abonnement~:} ACL (Galvanize/Diligent),
        IDEA (CaseWare), SAP GRC, Oracle GRC~; coût élevé, support complet~;
  \item \textbf{Outils open source ou gratuits~:} PSPP (analyse statistique),
        LibreOffice Calc avec macros, Python/pandas pour l'analyse de données~;
        coût nul, courbe d'apprentissage variable~;
  \item \textbf{Outils académiques et institutionnels~:} certains éditeurs
        (CaseWare IDEA) proposent des licences académiques ou des versions
        allégées adaptées aux institutions publiques.
\end{itemize}

\subsection*{Critère~5. Capacité TAAO (Techniques d'Audit Assistées par
Ordinateur)}

Un outil TAAO de qualité doit permettre~: l'importation de fichiers dans
des formats variés (CSV, Excel, XML, formats ERP propriétaires)~; la
réalisation des tests statistiques courants (stratification, échantillonnage,
détection de doublons, calcul de Benford)~; et la traçabilité complète
des étapes de traitement pour garantir la reproductibilité des tests.

\subsection*{Critère~6. Adaptation au contexte africain et institutionnel}

Pour les institutions d'audit public (CONSUPE, IGE, cours des comptes
africaines), trois contraintes additionnelles s'imposent~:
\begin{itemize}
  \item \textbf{Fonctionnement hors-ligne}~: les environnements audités
        peuvent ne pas avoir de connexion Internet stable~;
  \item \textbf{Formats locaux}~: les ERP utilisés (SIGIPES, IPPIS, IFMIS)
        exportent dans des formats spécifiques~;
  \item \textbf{Contraintes budgétaires}~: le coût par siège d'une licence
        SAP GRC ou ACL Enterprise peut être prohibitif pour une institution
        publique africaine.
\end{itemize}

% ------------------------------------------------------------------------------
\section{Classification fonctionnelle des outils}
\label{sec:classification-outils}
% ------------------------------------------------------------------------------

\begin{table}[H]
\centering
\small
\renewcommand{\arraystretch}{1.5}
\caption{Classification des outils d'audit par famille fonctionnelle}
\label{tab:classification-outils}
\begin{tabular}{|p{3cm}|p{5cm}|p{5.5cm}|}
  \hline
  \rowcolor{gray!20}
  \textbf{Famille} & \textbf{Description} & \textbf{Outils représentatifs} \\ \hline

  \textbf{F1. Analyse de données (TAAO)} &
  Extraction, interrogation, tests statistiques sur populations entières de
  transactions &
  ACL Desktop (Galvanize/Diligent), IDEA (CaseWare), Monarch (Datawatch),
  Python/pandas, R \\ \hline

  \textbf{F2. Gestion de mission d'audit} &
  Planification, papiers de travail, suivi des recommandations,
  reporting, archivage de la mission &
  ACL Workpapers, CaseWare Working Papers, TeamMate+, Protiviti PGP,
  RVR (Devoteam), Oxial Blue Audit \\ \hline

  \textbf{F3. Plateforme GRC intégrée} &
  Gouvernance, risques, conformité, contrôle interne, audit, plans de
  continuité~; plateforme unifiée multi-modules &
  SAP GRC~10, Oracle GRC, BWise (Nasdaq), eFront GRC, Enablon, Cura GRC,
  Cogis Software, Archer (RSA/EMC) \\ \hline

  \textbf{F4. Contrôle et audit continus} &
  Surveillance automatisée des transactions et des accès en temps réel
  ou quasi-réel, alertes &
  ACL AuditExchange, GreenLight Audit Analytics, modules SAP EGRCC,
  Oracle EGRC Controls \\ \hline

  \textbf{F5. Extraction de données complexes} &
  Lecture et transformation de fichiers dans des formats difficiles~:
  PDF semi-structurés, rapports d'ERP, fichiers mainframe &
  Monarch (Datawatch), ACL Importer, ACL Acerno, Python (pdfplumber,
  openpyxl), scripts ETL \\ \hline
\end{tabular}
\end{table}

% ------------------------------------------------------------------------------
\section{Fiches synthétiques des principaux outils}
\label{sec:fiches-outils}
% ------------------------------------------------------------------------------

\subsection*{ACL / Galvanize (Diligent)}

\begin{description}
  \item[Famille(s)] F1 (analyse de données), F2 (gestion de mission),
        F4 (contrôle continu).
  \item[Description] ACL est la référence mondiale de l'analyse de données
        en audit. Sa force est la capacité à analyser des populations entières
        de transactions sans échantillonnage, avec une traçabilité complète
        de chaque étape. La plateforme AuditExchange étend les capacités
        vers le contrôle continu centralisé. Le module Workpapers gère les
        papiers de travail avec lien direct vers les données analysées.
  \item[Points forts] Couverture fonctionnelle étendue~; intégration
        native des données SAP via ACL Direct Link~; traçabilité des tests~;
        large base d'utilisateurs et de ressources pédagogiques.
  \item[Points faibles] Coût de licence élevé~; courbe d'apprentissage
        significative pour les fonctions avancées~; dépendance à la
        connectivité pour les fonctions cloud.
  \item[Recommandation d'usage] Outil de référence pour les équipes d'audit
        interne de grande taille, les Big Four et les institutions d'audit
        public bien dotées. Idéal pour les missions à fort volume de
        transactions (paie, achats, stocks).
\end{description}

\subsection*{IDEA. CaseWare}

\begin{description}
  \item[Famille(s)] F1 (analyse de données), F2 (gestion de mission).
  \item[Description] IDEA (\emph{Interactive Data Extraction and Analysis})
        est reconnu pour sa simplicité d'utilisation relative et sa capacité
        à importer une grande variété de formats de fichiers. Il couvre les
        fonctions TAAO essentielles~: stratification, doublons, ruptures de
        séquence, calcul de la Loi de Benford, échantillonnage, analyse
        d'ancienneté, comparaison de fichiers. CaseWare Working Papers est
        le module compagnon de gestion des papiers de travail.
  \item[Points forts] Interface plus accessible qu'ACL~; large palette
        d'importateurs de formats (Excel, CSV, PDF, formats ERP)~; licences
        académiques disponibles~; documentation francophone disponible~;
        fonctionnement hors-ligne natif.
  \item[Points faibles] Moins adapté au contrôle continu que ACL
        AuditExchange~; capacités de visualisation limitées comparées aux
        outils modernes.
  \item[Recommandation d'usage] \textbf{Outil recommandé en premier choix
        pour les institutions d'audit public africaines}~: rapport
        fonctionnalités/coût favorable, fonctionnement hors-ligne, licences
        académiques disponibles, et adapté aux missions de taille moyenne.
        Convient parfaitement à l'audit de systèmes comme SIGIPES pour
        l'analyse des anomalies de paie.
\end{description}

\subsection*{Monarch. Datawatch/Altair}

\begin{description}
  \item[Famille(s)] F5 (extraction de données complexes).
  \item[Description] Monarch est un outil de transformation de données
        spécialisé dans l'extraction depuis des sources non structurées~:
        rapports PDF, fichiers texte à largeur fixe, exports d'imprimante,
        formats mainframe. Il permet de transformer ces sources en tableaux
        structurés exportables vers Excel, Access ou un outil d'analyse.
  \item[Points forts] Capacité d'extraction depuis des formats impossibles
        à lire directement par ACL ou IDEA~; aucune compétence de
        programmation requise~; interface visuelle intuitive.
  \item[Points faibles] Outil de transformation, pas d'analyse~; doit être
        utilisé en amont d'IDEA ou ACL.
  \item[Recommandation d'usage] Indispensable lorsque les données à analyser
        sont piégées dans des rapports PDF ou des formats d'impression. Outil
        de préparation des données, à coupler avec IDEA ou ACL pour les tests.
\end{description}

\subsection*{SAP GRC~10}

\begin{description}
  \item[Famille(s)] F3 (plateforme GRC), F4 (contrôle continu).
  \item[Description] SAP GRC est une suite intégrée de modules~: Access
        Control (gestion des risques liés aux autorisations et séparation
        des tâches dans SAP et non-SAP), Risk Management (gestion des
        risques d'entreprise), Process Control (contrôle interne et
        conformité), et Audit Management (planification et réalisation
        des missions).
  \item[Points forts] Intégration native avec SAP~: analyse des
        séparations de tâches (SoD) directement dans les données SAP~;
        contrôle continu des transactions~; référentiel de risques et
        de contrôles intégré.
  \item[Points faibles] Réservé aux environnements SAP~; coût de licence
        très élevé~; implémentation longue et coûteuse~; nécessite des
        compétences SAP spécialisées.
  \item[Recommandation d'usage] Pertinent uniquement dans les organisations
        dont le SI repose sur SAP. Inadapté aux institutions publiques
        africaines dont les SI sont généralement non-SAP.
\end{description}

\subsection*{BWise (Nasdaq)}

\begin{description}
  \item[Famille(s)] F3 (plateforme GRC intégrée).
  \item[Description] BWise est une plateforme GRC standard qui couvre
        l'audit interne, le contrôle interne, la gestion des risques et la
        conformité. Elle propose une description structurée de l'organisation
        (processus, risques, contrôles) et automatise les campagnes d'évaluation,
        les plans d'action et le reporting.
  \item[Points forts] Couverture fonctionnelle GRC complète~; indépendance
        vis-à-vis de l'ERP sous-jacent~; interfaces configurables.
  \item[Points faibles] Coût de déploiement et de licence élevé~; adapté
        aux grandes organisations~; peu accessible pour les petites équipes.
  \item[Recommandation d'usage] Adapté aux grandes DSI ou aux départements
        d'audit interne cherchant à unifier leurs processus GRC sur une
        plateforme unique. Non recommandé pour les missions ponctuelles.
\end{description}

\subsection*{RVR. Devoteam}

\begin{description}
  \item[Famille(s)] F2 (gestion de mission), F3 (GRC).
  \item[Description] RVR est une plateforme GRC française intégrée couvrant
        le processus d'audit (plan, planification, réalisation, papiers de
        travail, rapport, suivi des recommandations), le contrôle interne
        et la gestion des risques. Elle intègre les référentiels COSO et AMF
        nativement et suit une méthodologie conforme aux normes IIA.
  \item[Points forts] Interface francophone~; conformité IIA native~;
        intégration référentiels COSO/AMF~; adapté aux équipes d'audit
        francophones.
  \item[Points faibles] Moins connu hors de France~; écosystème de
        partenaires plus restreint qu'ACL ou SAP.
  \item[Recommandation d'usage] \textbf{Option pertinente pour les
        institutions d'audit public francophones d'Afrique subsaharienne}
        (CONSUPE, IGF, Cour des Comptes)~: interface et documentation
        francophone, conformité aux normes IIA, tarification plus accessible
        que les solutions anglo-saxonnes.
\end{description}

\subsection*{Python / pandas / openpyxl (open source)}

\begin{description}
  \item[Famille(s)] F1 (analyse de données), F5 (extraction).
  \item[Description] Python, avec ses bibliothèques d'analyse (pandas,
        numpy), de visualisation (matplotlib, seaborn) et de traitement
        de fichiers (openpyxl, pdfplumber), constitue une alternative
        open source complète aux outils TAAO commerciaux. Un auditeur
        maîtrisant Python peut reproduire l'ensemble des tests d'IDEA ou
        d'ACL Desktop sans coût de licence.
  \item[Points forts] Coût nul~; flexibilité maximale~; capacité à traiter
        des volumes très importants~; intégration avec les outils de
        visualisation~; reproductibilité totale (scripts versionnables sous
        Git).
  \item[Points faibles] Courbe d'apprentissage significative~; absence
        de traçabilité automatique des tests~; nécessite de construire
        soi-même les fonctions d'importation et de test~; non certifié
        par les organismes professionnels.
  \item[Recommandation d'usage] \textbf{Recommandé pour les auditeurs
        formés en génie informatique} (profil ENSP) qui disposent déjà
        de compétences en programmation~: coût nul, puissance maximale,
        et compétence directement valorisable dans les missions d'audit
        technique. À coupler avec un outil de gestion de mission (papiers
        de travail) pour une couverture complète.
\end{description}

% ------------------------------------------------------------------------------
\section{Arbre de décision pour le choix d'un outil}
\label{sec:arbre-decision-outils}
% ------------------------------------------------------------------------------

L'arbre ci-dessous guide le choix d'un outil selon quatre questions
successives. Il est conçu pour être opérationnel en contexte africain
et institutionnel.

\begin{enumerate}

  \item \textbf{La mission nécessite-t-elle principalement l'analyse de
        fichiers de données (transactions, accès, paie)~?}
  \begin{itemize}
    \item \textbf{Oui} → Aller à la question~2.
    \item \textbf{Non} (mission de gouvernance, de conformité, de
          gestion des risques) → Aller à la question~4.
  \end{itemize}

  \item \textbf{Disposez-vous d'un budget pour un outil commercial~?}
  \begin{itemize}
    \item \textbf{Oui, et l'équipe n'a pas de compétences en programmation}
          → \textbf{IDEA} (premier choix) ou \textbf{ACL Desktop}.
    \item \textbf{Oui, et l'organisation utilise SAP}
          → \textbf{SAP GRC + ACL Direct Link}.
    \item \textbf{Non, ou budget très limité} → Aller à la question~3.
  \end{itemize}

  \item \textbf{L'équipe dispose-t-elle de compétences en programmation
        Python~?}
  \begin{itemize}
    \item \textbf{Oui} → \textbf{Python/pandas/openpyxl}.
    \item \textbf{Non} → Demander une \textbf{licence académique IDEA}
          (CaseWare en propose pour les institutions d'enseignement)
          ou utiliser \textbf{Excel avec macros} pour les missions de
          taille limitée.
  \end{itemize}

  \item \textbf{La mission concerne-t-elle la gouvernance et la gestion
        des risques (pas l'analyse de transactions)~?}
  \begin{itemize}
    \item \textbf{Oui, contexte francophone, budget modéré}
          → \textbf{RVR} (Devoteam) ou tableur structuré (outil
          développé au chapitre~\ref{chap:outil-audit}).
    \item \textbf{Oui, grande organisation, budget élevé}
          → \textbf{BWise} ou \textbf{eFront GRC}.
    \item \textbf{Oui, données à extraire de formats non structurés
          (PDF, rapports)} → \textbf{Monarch} en amont, puis IDEA ou
          Python pour l'analyse.
  \end{itemize}

\end{enumerate}

\begin{boxalert}
\textbf{Recommandation pratique pour les institutions publiques africaines}

Pour les institutions d'audit public à budget contraint (CONSUPE, IGF,
corps de contrôle), la combinaison suivante offre le meilleur rapport
fonctionnalité/coût/autonomie~:

\begin{itemize}
  \item \textbf{IDEA} pour l'analyse de données (licence institutionnelle
        négociable)~;
  \item \textbf{L'outil Excel de décomposition hiérarchique}
        (chapitre~\ref{chap:outil-audit}) pour la gestion de la mission
        et la collecte des preuves~;
  \item \textbf{Python/pandas} pour les analystes ayant une formation
        en génie informatique, en particulier pour les analyses de fichiers
        volumineuses (paie, marchés publics).
\end{itemize}

Cette combinaison est entièrement fonctionnelle hors-ligne, ne dépend
d'aucun abonnement cloud, et peut être déployée sur des postes sans
connexion Internet~--- contrainte courante sur les terrains d'audit dans
les administrations publiques.
\end{boxalert}